\documentclass[10pt,a4paper]{amsart}
\usepackage[margin=19mm]{geometry}
\usepackage{amsmath,amssymb,mathtools}
\usepackage[scr=rsfs,cal=euler]{mathalfa}
\usepackage{xfrac}
\usepackage[unicode,hidelinks,bookmarks=false]{hyperref}
\newcommand{\ie}{i.e.\ }
\newcommand{\cf}{cf.\ }
\newcommand{\resp}{respectively\ }
\newcommand{\Subsection}{\subsection}
\providecommand{\nobreakdash}{\nobreak\hbox{-}}
\setlength{\emergencystretch}{2em}
\multlinegap0pt
\newcommand{\dd}{\mathrm{d}}
\theoremstyle{plain}
\newtheorem{Lemma}{Lemma}[section]
\newtheorem{Theorem}[Lemma]{Theorem}
\newtheorem{Proposition}[Lemma]{Proposition}
\newtheorem{Corollary}[Lemma]{Corollary}
\newtheorem{cor}[Lemma]{Corollary}
\newtheorem{prop}[Lemma]{Proposition}
\newtheorem{Conj}[Lemma]{Conjecture}
\newtheorem*{Corollary of Conjecture}{Corollary of Conjecture}
\newtheorem{qn}[Lemma]{Question}
\theoremstyle{definition}
\newtheorem{Def}[Lemma]{Definition}
\theoremstyle{remark}
\newtheorem{rem}[Lemma]{Remark}
\newtheorem{eg}[Lemma]{Example}
\newtheorem{ack}{Acknowledgments.}
\title[Conjugation Differential Invariants of $\mathrm{SL}_2(\mathbb{F}_q)$ on trace-free matrices]{Conjugation Differential Invariants of $\mathrm{SL}_2(\mathbb{F}_q)$ on trace-free matrices}
\author{Jonathan Elmer}
\date{Source: arXiv:2609.28195v2 (CC BY 4.0)}
\begin{document}
\maketitle

\noindent\emph{Source and licence.} The delivered work is the article shown in the title above, version arXiv:2609.28195v2 (CC BY 4.0), distributed under the Creative Commons Attribution 4.0 International licence, as recorded in the arXiv licence metadata for this exact version and shown on the abstract page saved with this item. The full licence notice, which carries the licence URL and the address of that abstract page, is the bundled file \texttt{licenses/arxiv-2609.28195-CC-BY-4.0.txt}.
\par\smallskip

\noindent\textit{Editorial scope.} Counted unit: the article's Corollary (source label oddmin) (source0.tex lines 248--253, source label \texttt{oddmin}): ``Suppose is odd. A minimal generating set for is given by f 1,f 2,f 3,f 4, f 1, f 2, f 3, ( f 2 f 3), ( f 3 f 1), ( f 1 f 2), -1 ( f 1), -1 ( f 2), -1 ( f 3), x 1 x 2 x 3. .''. It is delivered together with the article's complete original proof (source0.tex lines 686--701, one proof environment printed later in the article, detached from the statement by the article's own arrangement, 16 lines), copied line for line from the article's own text. Required context retained uncounted: none beyond the counted statement and proof. Editorial formatting only: an AMS \texttt{amsart} preamble that restates the article's own macro definitions and its own theorem declarations, and the theorem environments the retained lines use; the article's own class file is neither shipped nor input; the retained environments are renumbered sequentially in order of appearance, whereas the article prints them with its own numbering; the article's equation numbering is not reproduced and every cross-reference inside the retained text resolves to the wrapper's numbering. No citation occurs inside any retained line, so the wrapper carries no bibliography; All retained bodies are exact copies of the bundled \texttt{sources/211559/source0.tex}, so no omission receipt applies. No asset-inclusion command occurs inside any retained span and the package ships no image asset.


\section{The counted unit: Corollary (source label oddmin) and its complete original proof}

\begin{Corollary}\label{oddmin}
Suppose $q$ is odd. A minimal generating set for $\Omega^G$ is given by
\begin{align*}\{&f_1,f_2,f_3,f_4,\\ &\dd f_1, \dd f_2, \dd f_3, \phi(\dd f_2 \wedge \dd f_3), \phi(\dd f_3 \wedge \dd f_1), \phi(\dd f_1 \wedge \dd f_2),\\
&\phi^{-1}(\dd f_1), \phi^{-1}(\dd f_2), \phi^{-1}(\dd f_3), \\
& \dd x_1 \wedge \dd x_2 \wedge \dd x_3.\}.\end{align*}
\end{Corollary}

\begin{proof}[Proof of Corollary \ref{oddmin}]
It is clear that the set $\{f_1,f_2,f_3\} \cup S$ generates $\Omega^G$ and that $c_1,c_2,c_3$ and $f_4 \dd x_1 \wedge \dd x_2 \wedge \dd x_3$ can be omitted. It is clear that $\dd x_1 \wedge \dd x_2 \wedge \dd x_3$ cannot be omitted. Since $b_1, \ldots, b_6$ generate freely over $F[f_1,f_2,f_3]$ and each has $x_1$-degree strictly smaller than the $x_1$-degree of $f_4$, none of these can be ommited. It remains to show that none of $c_4,c_5, c_6$ can be omitted.

The lead terms of coefficients of these $2$-forms are given in the table below:

\vspace{1cm}
\begin{tabular}{|c||c|c|c|} \hline
 & $\dd x_2 \wedge \dd x_3$ &  $\dd x_3 \wedge \dd x_1$ &  $\dd x_1 \wedge \dd x_2$ \\ \hline \hline
$c_4$  & $x_1$ & $x_2$  & $x_3$\\ \hline
$c_5$  &  $x_1^q$ &  $x_2^q$ & $x_3^q$\\ \hline
$c_6$  &  $-\frac{q-1}{8}x_1^{\frac12 q(q-3)+2}x_2^{q-3}$ & $-\frac{q-1}{4}x_1^{\frac12 q(q-3)+1}x_2^{q-2}$  & $-\frac{q-1}{2}x_1^{\frac12 q(q-3)}x_2^{q-1}$\\ \hline
\end{tabular}
\vspace{1cm}

Again each has $x_1$-degree strictly smaller than that of $f_4$, so we see that $\{c_1,c_2, \ldots, c_6\}$ is a minimal generating set for $(\Omega^2)^G$ over $R^G$. It remains to show that none of $c_4,c_5,c_6$ lie in the $\mathcal{A}$-submodule of $\Omega^2$ generated by wedge products of $b_4,b_5, b_6$. But any such wedge product has $x_1$-degree at least $\frac12 q(q-1)$ which is greater than the $x_1$-degree of any of these.
 \end{proof}

\end{document}
